% Generated draft: requires mathematical and source review.
\documentclass[11pt]{article}
\usepackage[margin=25mm]{geometry}
\usepackage{fontspec}
\setmainfont{Latin Modern Roman}
\usepackage{amsmath,amssymb,mathtools,unicode-math}
\setmathfont{Latin Modern Math}
\usepackage{enumitem,needspace,xurl,hyperref}
\setlength{\parindent}{0pt}
\setlength{\parskip}{5pt}
\setlength{\emergencystretch}{3em}
\newcommand{\N}{\mathbb{N}}
\newcommand{\Z}{\mathbb{Z}}
\newcommand{\Q}{\mathbb{Q}}
\newcommand{\R}{\mathbb{R}}
\newcommand{\C}{\mathbb{C}}
\newcommand{\F}{\mathbb{F}}
\newcommand{\A}{\mathbb{A}}
\newcommand{\PP}{\mathbb{P}}
\newcommand{\E}{\mathbb{E}}
\newcommand{\eps}{\varepsilon}
\newcommand{\dd}{\,\mathrm{d}}
\begin{document}
\section*{UnsolvedMath ID: 30002777}
\textbf{Source problem code:} catalogue entry
\subsection{Definitions and mathematical statement}
Let \(k\) be a coefficient field, whose intended hypotheses are not specified in the supplied record. Let \(\Delta\) be a finite simplicial complex on vertex set \(V=\{1,\ldots,n\}\). Using the usual rank convention, put
\[
s=\max\{|F|:F\in\Delta\}=\dim\Delta+1.
\]
Its Stanley--Reisner ring over \(k\) is the standard graded ring
\[
R=k[\Delta]=k[x_1,\ldots,x_n]/I_\Delta,
\qquad
I_\Delta=\left(\prod_{j\in F}x_j:F\subseteq V,\ F\notin\Delta\right),
\qquad \deg x_j=1.
\]
The complex is Cohen--Macaulay over \(k\) if \(\operatorname{depth}R=\dim R=s\). It is doubly Cohen--Macaulay over \(k\) if it is Cohen--Macaulay over \(k\) and, for every vertex \(v\in V\), the deletion
\[
\Delta-v=\{F\in\Delta:v\notin F\}
\]
is Cohen--Macaulay over \(k\) and has the same dimension as \(\Delta\).

A linear system of parameters is a sequence \(\Theta=(\theta_1,\ldots,\theta_s)\) with \(\theta_j\in R_1\) such that
\[
A=R/(\theta_1,\ldots,\theta_s)R
\]
is finite dimensional over \(k\). Write \(A_i\) for its homogeneous component of degree \(i\).

For a fixed coefficient field \(k\), the supplied question has the following conditional formulation: does every finite simplicial complex \(\Delta\) that is doubly Cohen--Macaulay over \(k\) admit a linear system of parameters \(\Theta\) and a single linear form \(\omega\in R_1\) such that, simultaneously for every integer \(i\) with \(0\leq i\leq\lfloor s/2\rfloor\), multiplication by the indicated power of the image \(\bar\omega\in A_1\) is injective?
\[
\times\bar\omega^{\,s-2i}:A_i\longrightarrow A_{s-i},
\qquad a\longmapsto\bar\omega^{\,s-2i}a.
\]
The choices of \(\Theta\) and \(\omega\) may depend on \(\Delta\) and \(k\), but must work for all these degrees at once.

The coefficient-field quantifier remains unresolved: the record does not specify characteristic, infinitude, or whether scalar extension is allowed before choosing \(\Theta\) and \(\omega\). Thus the formulation above does not assert that every field is intended. The original wording refers to the injectivity conclusion of ``Theorem 4'' and uses \(\#S\), but neither that theorem nor the definition of \(S\) is supplied. The displayed formulation follows the supplied cleaned statement, which replaces \(\#S\) by the rank \(s\); the precise connection to the original theorem requires review.
\subsection{Short English statement}
For every doubly Cohen--Macaulay simplicial complex, can one choose a linear system of parameters and one linear form whose powers inject each degree in the lower half of the resulting quotient into the complementary degree? The intended coefficient-field hypotheses and the original meaning of \(\#S\) remain unspecified.
\subsection{Sources}
\begin{itemize}
\item Supplied local UnsolvedMath registry record, permanent ID \(30002777\), source problem code \texttt{catalogue entry}, title ``Lefschetz Injectivity for Doubly Cohen--Macaulay Complexes,'' in the Oberwolfach Reports Open Problems collection; proposed year recorded as 2015. No individual proposer is identified.
\item \emph{Geometric and Algebraic Combinatorics} (2015), Oberwolfach Reports, pp.~285--368, DOI 10.4171/owr/2015/5: \url{https://doi.org/10.4171/owr/2015/5}. The supplied literature-source label is ``Murai, Oberwolfach Report.'' The background describes these reports as records of workshops and their open-problem discussions.
\item Supplied literature reference labeled ``Juhnke--Kubitzke--Murai, primary paper'': \url{https://arxiv.org/abs/1503.06430}.
\item Literature assessment dated \texttt{2026-08-21}: supplied dated curation, not independently verified here. It labels the problem partially solved and reports that the Lefschetz-type conclusion is proved for rank-selected subcomplexes of balanced simplicial polytopes, while the assertion for all doubly Cohen--Macaulay complexes remains open. The record explicitly calls for independent verification of time-sensitive status and cited claims.
\item The registry labels the statement \texttt{corrected\_verified} and reports that the official Oberwolfach passage was checked directly and its question, definitions, and notation restored without strengthening the claim. This is supplied verification metadata; no independent source inspection or mathematical verification is asserted here.
\end{itemize}
\end{document}
